\chapter{Введение}
\label{ch:intro}

\section{Введение}

Команда netstat предоставляет информацию о различных сетевых структурах данных в
различных форматах в зависимости от указанных опций. В отличие от Wireshark, Netstat не
показывает абсолютно все перехваченные пакеты в реальном времени, а отображает статус
различных соединений (TCP/UDP). \\

Сокет – это сочетания IP-адреса источника с номером порта источника и IP-адреса
назначения с номером порта назначения. Такая связка «IP-адрес:Порт» позволяет однозначно
определить конечную точку доставки сообщения, потому что оно будет адресовано не просто
какому-то устройству, а конкретному приложению на данном устройстве.


\section{Цель}  

\begin{enumerate}
    \item Выполнить установку и вывести результаты команды Netstat;
    \item Изучить и использовать правильный синтаксис команды Netstat;
    \item Проанализировать и определить состояние сокетов.
\end{enumerate}

\endinput 